%%=============================================================================
%% Inleiding
%%=============================================================================

\chapter{\IfLanguageName{dutch}{Inleiding}{Introduction}}%
\label{ch:inleiding}

24flow is een bedrijf dat een modulair platform maakt voor werkvloerbeheer en productieplanning om onder andere de traceerbaarheid en zichtbaarheid te verbeteren bij make-to-order fabrikanten. Bedrijven die focussen op maatwerk en lage volumes vormen hierbij de doelgroep. Het 24flow-platform zorgt voor beter begrip en overzicht van de productieprocessen en biedt hiermee oplossingen aan om de doorlooptijd van deze processen te verminderen en de kosten te drukken.
\\
Dit platform is volledig gebouwd binnen Salesforce, een Customer Relationship Management (CRM) platform dat gebruikt wordt om bedrijven hun klantgegevens, operaties, verkoop en andere processen te laten beheren. 

\section{\IfLanguageName{dutch}{Probleemstelling}{Problem Statement}}%
\label{sec:probleemstelling}

Momenteel hebben (nieuwe) gebruikers op het platform moeite om alle functionaliteiten zelfstandig te gebruiken en te begrijpen. Er moet vaak handmatig uitgelegd of getoond worden hoe alles werkt, en dit is bij elke vernieuwing of update ook nog eens extra het geval. Er zijn tooltips en tutorials beschikbaar binnen het platform, maar deze zijn beperkt of onduidelijk. Al deze zaken kosten veel tijd en geld, en kunnen bovendien voor frustraties zorgen bij de klanten.
\\

\section{\IfLanguageName{dutch}{Onderzoeksvraag}{Research question}}%
\label{sec:onderzoeksvraag}

Hoe kan het onboardingproces en de in-app guidance van (nieuwe) gebruikers op het 24flow-platform, gebouwd binnen het Salesforce-ecosysteem, worden verbeterd, en hoe kunnen bestaande Salesforce-functionaliteiten hierbij helpen?
\\

\section{\IfLanguageName{dutch}{Onderzoeksdoelstelling}{Research objective}}%
\label{sec:onderzoeksdoelstelling}

Om het probleem beter te begrijpen, moeten we enke ons afvragen hoe het huidige onboardingproces in elkaar zit en op welke manieren gebruikers begeleid worden binnen het platform, welke bestaande frustraties of gebreken huidige klanten, gebruikers en ontwikkelaars ondervinden en hoe andere vergelijkbare systemen omgaan met onboarding en in-app guidance.
\\

Door een antwoord te vinden op deze vragen, zal dit onderzoek pogen tot een eindresultaat te komen dat als proof of concept een in-app guidance systeem met walkthroughs ontwikkelt binnen het 24flow-platform. Om hiertoe te komen kijken we welke Salesforce (Apex) mogelijkheden er bestaan voor onboarding en guidance, welke componenten en processen het meeste nood hebben aan user guidance, en hoe andere vergelijkbare systemen omgaan met onboarding.
\\

Om het probleem beter te begrijpen, is dit onderzoek opgebouwd rond een aantal gerichte deelvragen. Binnen het \textbf{probleemdomein} staan de volgende vragen centraal:

\begin{enumerate}
	\item Hoe zit het huidige onboardingproces in elkaar en op welke manieren worden gebruikers momenteel begeleid binnen het platform?
	\item Welke bestaande frustraties of gebreken ondervinden gebruikers en ontwikkelaars bij de onboarding?
	\item Hoe gaan andere, vergelijkbare systemen om met onboarding en in-app guidance?
\end{enumerate}

Door een antwoord te vinden op deze vragen, zal dit onderzoek pogen tot een eindresultaat te komen dat als proof of concept een in-app guidance systeem met walkthroughs ontwikkelt binnen het 24flow-platform. Om dit proof of concept succesvol te realiseren, worden de volgende deelvragen binnen het \textbf{oplossingsdomein} beantwoord:

\begin{enumerate}
	\item Welke Salesforce (Apex) mogelijkheden bestaan er voor het opzetten van onboarding en user guidance?
	\item Welke componenten en processen binnen het 24flow platform hebben het meeste nood aan user guidance?
\end{enumerate}
\section{\IfLanguageName{dutch}{Opzet van deze bachelorproef}{Structure of this bachelor thesis}}%
\label{sec:opzet-bachelorproef}

De rest van deze bachelorproef is opgebouwd volgens de opeenvolgende fasen van het onderzoek:
\\

In Hoofdstuk~\ref{ch:stand-van-zaken} wordt de stand van zaken toegelicht. Dit hoofdstuk bevat een uitgebreide literatuurstudie naar software-onboarding en een diepgaande technologische analyse van de native Salesforce-functionaliteiten.
\\

In Hoofdstuk~\ref{ch:methodologie} wordt de methodologie uiteengezet en worden de gebruikte onderzoekstechnieken besproken om een gestructureerd antwoord te kunnen formuleren op de onderzoeksvragen.
\\

In Hoofdstuk~\ref{ch:requirements-ontwerp} worden de functionele en niet-functionele vereisten voor het proof of concept vastgelegd via de MoSCoW-methode en visueel uitgewerkt in use cases en wireframes.
\\

In Hoofdstuk~\ref{ch:realisatie} wordt de daadwerkelijke technische realisatie en de architectuur van de applicatie (zowel de front-end als de backend-controller) gedetailleerd toegelicht.
\\

In Hoofdstuk~\ref{ch:evaluatie} wordt het ontwikkelde prototype kritisch geëvalueerd aan de hand van de requirements en de technische platformbeperkingen.
\\

In Hoofdstuk~\ref{ch:conclusie}, ten slotte, wordt het algemene besluit gevormd en een definitief antwoord geformuleerd op de onderzoeksvragen, inclusief een blik op toekomstige ontwikkelingen binnen dit domein.